\section[The unbounded-b source with xb infinitesimal]{The unbounded-$b$ source with $xb$ infinitesimal}
\label{sec:c2-unbounded-b-small-B}

\begin{theorem}[The remaining unbounded-state rational C2 source]
Use the exact original rational C2 source.  Suppose $x>0$ is
infinitesimal, $k$ is a positive bounded unit, $b$ is positive
unbounded, $B=xb$ is infinitesimal, $0<r<1$, $0<h<b$, and all
original positive-chart inequalities hold.  Assume $n\geq0$ and
physicality of the source.  Then the full original parent has tensor
SOS for every independent real or complex projective state and
every independent projective pair of measurement angles.
\end{theorem}

The conclusion includes $n=0$, arbitrarily large $n$, arbitrary
complementary openings, and every value of the pure coefficient
ratio $B/\sigma$, where $\sigma=1/b$.  In particular, the actual
contact ratio $W/Z=k^{-1}\sqrt{B/\sigma}$ need not be bounded.
This proof uses the homogeneous scalar lemma in
Section~\ref{sec:unbounded-b-homogeneous-scalar}; it restores both
full state normals and exhausts the other testing charts.

\paragraph{Exact source scales.}
Put
\[
 \sigma=1/b,\quad x=\sigma B,\quad \eta=h/b,\quad
 y=B\eta,\quad z=B(1-\eta),\quad
 A_c=1+\sigma Br,\quad S_c=A_c+z.
\]
The positive chart gives $0<\eta<1$; both $y,z$ are infinitesimal.
Let $\widehat\Delta=\sigma\Delta$,
$\widehat N=\sigma^2 B N_{\rm src}$,
$\widehat g=\sigma g=1+\sigma(r-1+B\eta)$ and
$j=A_c-z/k^2$.  The cancelled exact formulas are
\begin{align*}
 nr&=\frac{(1+\sigma r)A_cS_c-r\widehat\Delta}
                       {2\widehat\Delta},\\
 m&=\frac{\sigma^2(1+k^2)(1-r)A_cS_c}
                       {2\widehat g\widehat\Delta},\\
 C&=\frac{A_c\widehat N}{2z\widehat\Delta},\qquad
 A=\frac{k^2A_c(1+\sigma r)S_c}{2y\widehat\Delta},\\
 D&=\frac{B(1+\sigma B)(1+\sigma r)(1-\sigma k^2r)S_c}
                       {2k^2j\widehat\Delta},\\
 \alpha&=\frac{B(1+\sigma r)\widehat N}
                       {2(1-y)\widehat\Delta}.
                                                               \tag{UB1}
\end{align*}
Here $\widehat\Delta,\widehat N,\widehat g,j,1-y$ are bounded
positive units.  Write $K=k_0>0$ and $D_K=1+K^2$.  Their relevant
residues imply
\begin{gather*}
 (nr)_0=\tfrac12(D_K^{-1}-r_0),\quad
 (m/\sigma^2)_0=(1-r_0)/2,\\
 (Cz)_0=(Ay)_0=K^2/(2D_K),\quad
 (D/B)_0=1/(2K^2D_K),\quad
 (\alpha/B)_0=K^2/(2D_K).                            \tag{UB2}
\end{gather*}
Since $n\geq0$, $r_0\leq D_K^{-1}<1$, so $1-r$ is a unit.
Set $N=4n+2\geq2$.  Exactly,
\[
 Nr=\frac{2(1+\sigma r)A_cS_c}{\widehat\Delta},      \tag{UB3}
\]
which is a bounded positive unit.  Thus $nr,nr^2$ are bounded even
when $n$ is unbounded, while $1/N$ and $1/\sqrt N$ are bounded.

All opening cancellations needed below are uniform:
\[
 BC(1-\eta)=Cz,\quad BA\eta=Ay,\quad
 BC(1-\eta)^2=Cz(1-\eta),\quad BA\eta^2=Ay\eta.
                                                               \tag{UB4}
\]
The quantity $H_B=B(C+A)$ is bounded below by a positive real;
its reciprocal and reciprocal square root are bounded.  No
assumption that $\eta$ or $1-\eta$ is a unit is made.

\paragraph{The whole finite $u=0$ axis.}
Use the original parent and the exact-contact graph
\[
 X_g=-ruZ,\qquad Y_g=\eta uW,
 \qquad E_X=\sqrt N(X+ruZ),\quad
 E_Y=\sqrt{C+A}(Y-\eta uW).                          \tag{UB5}
\]
The actual contact is
\[
 u=\sqrt{\sigma B},\quad v=k,\quad
 W=k^{-1}\sqrt{B/\sigma},\quad Z=1.
\]
Let $\xi=u^2$ and $f=F|_g$.  Exact polynomial coquotients give
\[
 f=a(\xi,v)W^2+2u d(\xi,v)WZ+c(\xi,v)Z^2,
 \quad a=\frac\xi B Q+\sigma^2a_0,\quad c=B e+\xi b.
                                                               \tag{UB6}
\]
Every coefficient of $Q,a_0,d,e,b$ is bounded by (UB3)--(UB4).
For example $a_0=4(m/\sigma^2)v^2$,
$e=4(D/B)v^2+4\alpha/B$, and
$Q=B(a-a|_{u=0})/u^2$.  These are monic polynomial divisions.
The full leading functions at $u=0$ are
\begin{align*}
 q(v)&=2K^2(1+v^2)/D_K,& d_0(v)&=-2v,\\
 e_0(v)&=2(v^2/K^2+K^2)/D_K,&
 (a_0)_0(v)&=2(1-r_0)v^2,\\
 b_0(v)&=2r_0(1+v^2)/D_K-4r_0+2.                  \tag{UB7}
\end{align*}

For remote tests use
\[
 P=uW/\sqrt B,\quad B_u=uZ,\quad
 V=\sigma vW,\quad E=\sqrt B Z.
\]
Then (UB6) has the exact bounded Gram
\[
 QP^2+b B_u^2+(a_0/v^2)V^2+eE^2+2dPE,             \tag{UB8}
\]
where $a_0/v^2$ is a polynomial coquotient.  The $(P,E)$ residue
block is strictly positive whenever $v_0^2\ne K^2$, since
\[
 q(v)e_0(v)-4v^2=4(v^2-K^2)^2/D_K^2.               \tag{UB9}
\]
This is a global identity for every finite residue $v$, not a
local-contact positivity inference.

\paragraph{Both complete state normals.}
The exact quadratic normal expansion is
\[
 F=f+E_X L_X+E_Y L_Y+H_{XX}E_X^2
                       +2H_{XY}E_XE_Y+H_{YY}E_Y^2.
                                                               \tag{UB10}
\]
Its entire hard block is bounded with residue
$\operatorname{diag}(D_{v,0},4D_{v,0})$.
Indeed the $X$ pivot is $ND_uD_v+O(1)$ before division by $N$,
with the remaining terms either infinitesimal in $u$ or of order
$m$; the $Y$ pivot from the strong packets is exactly
$4D_v(1+[A/(C+A)]u^2)$.  All mixed entries have zero residue and
contain no $C,A$ contribution.

The complete unnormalized $X$ derivative has the form
\[
 F_X|_g=uZ A_X+W(\xi B_1+\sigma^2 B_0),             \tag{UB11}
\]
with bounded polynomial coefficients, where
\begin{align*}
 B_0&=-8(m/\sigma^2)v(1+\eta\xi),\\
 B_1&=v(8D+4)(1-\eta),\\
 A_X&=-8nrD_uD_v+8m(1-r)
       -(8D+4)v^2(1+r\xi)+4D_v(1-r).
\end{align*}
Thus in the remote roots its $P,V$ coefficients have zero residue;
its only possible nonzero residue coupling is to $B_u$.
The exact residue Schur expression for the $(B_u,E_X)$ block is
\[
 (N r^2D_v-4r+2)
       -\frac{(4-2NrD_v)^2}{4ND_v}
                 =2-\frac4{ND_v}.                  \tag{UB12}
\]
It is strictly positive for every finite nonzero residue $v$,
because $N\geq2$.  At $v_0=0$ it is strict if $n$ is unbounded
or has positive residue.  The sole zero case $n_0=v_0=0$ is
treated below.

For the other complete first normal,
\[
 L_Y=\frac1{\sqrt{H_B}}
 \{P(B_2+B e_w)+B_u(\sqrt B\,u e_z)+E e_*\},       \tag{UB13}
\]
where $e_w,e_z$ are bounded,
\[
 B_2=8D_v\{(BA)\eta D_u-(BC)(1-\eta)\},\qquad
 e_*=8v(D-\alpha).
\]
All coefficients have zero residue, by (UB4), $B\to0$, and
$u\to0$.  This verifies both full normals and proves the full
parent at every remote finite-axis test covered by (UB12).

\paragraph{Matched residues without any state-ratio restriction.}
For $v_0=K$, apply the homogeneous scalar lemma to (UB6).  Its
contact and full stationarity hypotheses hold exactly because the
graph (UB5) passes through the original stationary contact.
Use the homogeneous roots
\[
 R=\frac{uW-(B/k)Z}{\sqrt B},\quad
 S=uZ-\sigma kW,\quad
 U=\sqrt B(v-k)Z,\quad T=u(v-k)Z.                  \tag{UB14}
\]
They satisfy
\[
 (\xi-x)W=u\sqrt B R+(B/k)S,\qquad
 (\xi-x)Z=uS+\sigma k\sqrt B R.                    \tag{UB15}
\]
The lemma constructs an exact bounded scalar Gram, valid for
every $W,Z$ simultaneously.  Its $(R,U)$ residue is
\[
 2K^2R^2+\frac{4(K^2-1)}{D_K}RU+\frac2{K^2}U^2,
\]
whose Schur margin is $8/D_K^2>0$.  Its radial $S^2$ residue is
$s_0=2(1-r_0)>0$; its only other possible residue coupling is
$\ell ST$, with a bounded real residue $\ell$.

To lift the full $X$ normal, set
\[
 L=A_X+(B/k)B_1+(\sigma/k)B_0.
\]
Equations (UB11) and (UB15) give exactly
\[
 F_X|_g=u\sqrt B B_1R-(\sigma/k)B_0S+uZL.
\]
At the source contact $L(x,k)=0$, because $F_X=0$ there.
Write its bounded monic differences as
$L=(\xi-x)L_\xi+(v-k)L_v$.  Then
\[
 F_X|_g=
 u\sqrt B(B_1+\sigma kL_\xi)R
 +\{\xi L_\xi-(\sigma/k)B_0\}S+L_vT.             \tag{UB16}
\]
After division by $\sqrt N$, all coefficients are bounded.
The $R,S$ coefficients have zero residue; the $T$ coefficient
may have a finite residue $\gamma$ (it is
$-8K/(D_K\sqrt{N_0})$ when $N$ is bounded).
This nonzero full-normal coupling is retained, not silently
discarded.

For (UB13), put $A_Y=B_2+B e_w$ and
$L_Y'=A_Y/k+\xi e_z+e_*$.  Its numerator is exactly
$RA_Y+\sqrt B ZL_Y'$.  Exact normal contact gives
$L_Y'(x,k)=0$.  Monic differences and (UB15) yield bounded
coefficients on $R,S,U$, all of zero residue: the $R,S$ corrections
carry factors $\sigma B$ or $u\sqrt B$, and the $U$ coefficient
is the angular coquotient of an identically zero leading function.
Division by $\sqrt{H_B}$ preserves these bounds.

Finally add $\delta$ times the exact zero
\[
       T^2-\sigma U^2-ST(v-k)-\sigma kRU(v-k)=0.    \tag{UB17}
\]
Choose a real $\delta>\ell^2/(4s_0)+\gamma^2/(4D_K)$.
This supplies a $T^2$ residue $\delta$ at zero residue cost to
the $U^2$ coefficient, since $\sigma$ is infinitesimal.  The block
on $(S,E_X,T)$ has strictly positive Schur complement; the $(R,U)$
block and $E_Y$ diagonal are already strict.  Thus the complete
six-root Gram on $(R,S,U,T,E_X,E_Y)$ is strictly positive in
residue.  The calculation is homogeneous in all four state
coordinates and needs no bound on $W/Z$, no comparison of $B$ with
$\sigma$, and no state endpoint exclusion.  The exact simultaneous
sign copy treats $v_0=-K$.

\paragraph{The double-zero test, including literal $n=0$.}
Suppose now $n_0=v_0=0$.  Use the pure coefficient quantities
\[
 H_c=4n+4m+2,\quad \zeta=(4m+2)/H_c,\quad
 \xi_c=4m/H_c,\quad \eta_c=C/(C+A),
\]
and graph $X_g=\xi_cvW-\zeta uZ$, $Y_g=\eta_cuW$.
Normalize its $X$ normal by $\sqrt{H_c}$ and its $Y$ normal by
$\sqrt{C+A}$.  Put
\[
 \epsilon=4\alpha+4n(4m+2)/H_c>0.
\]
Then $\epsilon$ is infinitesimal, $n/\epsilon$ and $B/\epsilon$
are bounded, and the harmonic product $BCA/(C+A)$ is a unit.
Use the five scalar roots
\[
 P=uW/\sqrt B,\quad V=\sigma vW,\quad E=\sqrt B Z,
 \quad B_\epsilon=\sqrt\epsilon\,uZ,\quad T_0=uvZ.
\]
The exact inherited identity, with its original coefficients,
is
\[
 f_{zz}=4\alpha+\epsilon u^2+4n\zeta^2u^4
                       +4Dv^2+u^2v^2 B_b,
 \qquad (B_b)_0=2.
\]
The $W^2$ coefficient at $u=0$ is exactly
\[
 \frac{4m(2n+1)}{2n+2m+1}v^2
              +(4n+2)\xi_c^2v^4.
\]
Therefore its coefficient on $V^2$ has positive residue
$4(m/\sigma^2)_0=2(1-r_0)$.  The scalar residue on the five roots
is strictly diagonal, with entries
\[
 \left(\frac{2K^2}{D_K},\ 2(1-r_0),\
        \frac{2K^2}{D_K},\ 1,\ 2\right).
\]
The cross term is allocated using $uvWZ=vPE$, so it has zero residue.

Both full normals are restored, including when $n=0$ exactly.
At $v=0$ the $X$ derivative is $-8n\zeta u^3Z$, whose coefficient
on $B_\epsilon$ is infinitesimal since $n/\epsilon$ is bounded.
Its remaining $Z$ coefficient is $uv^2$ times a bounded function,
so it has zero residue on $T_0$.  Its $W$ coefficient is a sum of
$u^2vW$ times a bounded function and $mvW$ times a bounded
function divisible by $v^2$ at $u=0$.  Allocate the former to
$P$ with coefficient $u\sqrt B\,v$, and the latter to $V$ with
coefficient $m/\sigma$; both have zero residue.  In particular one
must not divide the former by $\sigma$.
The $Y$ derivative is the inherited harmonic-opening identity:
$B_2=8D_v(BA)\eta_cu^2$, with bounded residual coquotients.
Its coefficients on $P,T_0,E$ all have zero residue.  The hard
residue is $\operatorname{diag}(1,4)$.  This proves the full-parent
corner without deleting a state coordinate.

\paragraph{All remaining projective tests.}
For the other commuting charts use the four original outer-core
identities with $\rho$ replaced by $\sqrt B$, and cap $C,A$ at
$s/B$, where $s=\min(1,BC,BA)$ is a positive bounded unit.
Retain both original nonnegative excess packets and the entire
original $n$ packet.  No positive lower bound on $m$ is needed:
the first two outer pivots are $4m_0+2=2$, while the other two
charts retain the whole $m$ packet and do not use it for strictness.
The reciprocal-$v$ chart uses $(D/B)_0>0$ exactly as before.
Thus the finite unit $u$ axis, both reciprocal axes, all double
corners, and literal projective endpoints have strict bounded
Grams.  The inherited physical noncommuting-residue theorem covers
all remaining testing pairs.

The finite-axis certificates above are homogeneous in the state,
including the previously problematic unbounded actual state
ratio.  Splitting a complex state into its real and imaginary
parts proves the complex assertion without an additional chart.
All inverses and radicals used in coefficients are pure coefficient
or pure testing operations; every displayed mixed quotient is
explicit monic polynomial division.  The roots may be unbounded.
The strict bounded tensor Gram lemma and restoration of the retained
packets complete the all-tests proof.
